\subsubsection{Pooling Size and Stride}\label{sec:pooling-size-and-stride}

In the previous section, we have seen that pooling layers are used to reduce the spatial dimensions of the input feature maps. But, \emph{\textbf{how does the pooling window move across the feature map?}} The movement of the pooling window is determined by two parameters: the \textbf{pooling size} and the \textbf{stride}. Typically, the \hl{stride is assumed equal to the pooling size.}

\highspace
\begin{flushleft}
    \textcolor{Green3}{\faIcon{th-large} \text{Pooling size}}
\end{flushleft}
The \textbf{pooling size} specifies the \hl{spatial neighborhood over which we compute the maximum.}

\highspace
For example, with $K = 2 \times 2$, each output activation summarizes a $2 \times 2$ region of the input feature map. And as we have seen, this is done \textbf{independently for every channel} of the input feature map.

\highspace
\begin{flushleft}
    \textcolor{Green3}{\faIcon{arrows-alt-h} \text{Stride}}
\end{flushleft}
The \textbf{stride} $S$ (\autopageref{sec:stride}) tells us \hl{how far the pooling window moves between two consecutive operations.} The most common configuration is to set the stride equal to the pooling size, i.e., $S = K$.

\highspace
For example:
\begin{equation*}
    K = 2 \times 2, \quad S = 2 \times 2
\end{equation*}
The first $2 \times 2$ region is processed, then the window moves $2$ \textbf{positions}, so the regions do not overlap. If we have one feature map of size $8 \times 8$, and apply a $2 \times 2$ max pooling with stride $2$, along the width, 8 positions are divided into groups of 2:
\begin{equation*}
    \boxed{1 \; 2} \quad \boxed{3 \; 4} \quad \boxed{5 \; 6} \quad \boxed{7 \; 8}
\end{equation*}
From each pair/group, pooling produces one output position. Therefore, the output width is $4$. The same happens along the height, so the output height is also $4$. The output feature map has size $4 \times 4$. Consequently, the \textbf{total number of spatial activations} becomes:
\begin{equation*}
    \text{Reduction Width} \times \text{Reduction Height} = \dfrac{1}{2} \times \dfrac{1}{2} = \dfrac{1}{4}
\end{equation*}
So a $2 \times 2$ max pooling with stride 2 \textbf{reduces the spatial size to 25\%}, or equivalently \hl{discards 75\% of the spatial samples.} However, the number of channels remains the same.

\highspace
\begin{flushleft}
    \textcolor{Red2}{\faIcon{question-circle} \textbf{What if stride is 1?}}
\end{flushleft}
Suppose we have a $2 \times 2$ max pooling with stride $1$:
\begin{equation*}
    K = 2 \times 2, \quad S = 1
\end{equation*}
Now the pooling windows \textbf{overlap} because the window moves only one position at a time. So we're computing a \textbf{local maximum around essentially every spatial position} of the input feature map (see \autoref{fig:pooling-stride-comparison}). So, stride $1$ can be used when pooling is intended primarily as a \textbf{local nonlinear maximum operation}, rather than primarily for spatial reduction.

\highspace
\begin{examplebox}[: Pooling Size and Stride]
    Suppose a convolutional filter produced this $4 \times 4$ feature map:
    \begin{equation*}
        X = \begin{bmatrix}
            1 & 2 & 1 & 0 \\
            3 & \mathbf{9} & 2 & 1 \\
            1 & 2 & 4 & 2 \\
            0 & 1 & 2 & 3
        \end{bmatrix}
    \end{equation*}
    The bold value $9$ means that \hl{the convolutional filter detected its feature very strongly at that position.} Now, we apply a $2 \times 2$ max pooling:
    \begin{itemize}
        \item Stride $S = 2$ (non-overlapping), we move the window \textbf{two positions at a time}:
        \begin{gather*}
            \begin{bmatrix}
                1 & 2 \\
                3 & \mathbf{9}
            \end{bmatrix} \to 9, \quad
            \begin{bmatrix}
                1 & 0 \\
                2 & 1
            \end{bmatrix} \to 2, \\
            \begin{bmatrix}
                1 & 2 \\
                0 & 1
            \end{bmatrix} \to 2, \quad
            \begin{bmatrix}
                4 & 2 \\
                2 & 3
            \end{bmatrix} \to 4
        \end{gather*}
        Therefore:
        \begin{equation*}
            \begin{bmatrix}
                1 & 2 & 1 & 0 \\
                3 & 9 & 2 & 1 \\
                1 & 2 & 4 & 2 \\
                0 & 1 & 2 & 3
            \end{bmatrix}
            \xrightarrow[\text{stride 2}]{2 \times 2 \text{ max pool}}
            \begin{bmatrix}
                9 & 2 \\
                2 & 4
            \end{bmatrix}
        \end{equation*}
        We went from $4 \times 4$ to $2 \times 2$, discarding $75\%$ of the spatial samples. So here pooling is clearly being used for \textbf{downsampling}.


        \item Stride $S = 1$ (overlapping), we move the window \textbf{one position at a time}. First window:
        \begin{equation*}
            \begin{bmatrix}
                1 & 2 \\
                3 & \mathbf{9}
            \end{bmatrix} \to 9
        \end{equation*}
        Move one position to the right:
        \begin{equation*}
            \begin{bmatrix}
                2 & 1 \\
                \mathbf{9} & 2
            \end{bmatrix} \to 9
        \end{equation*}
        Move one position right again:
        \begin{equation*}
            \begin{bmatrix}
                1 & 0 \\
                2 & 1
            \end{bmatrix} \to 2
        \end{equation*}
        Then we repeat the same process for the next rows, moving one position at a time. The final output is:
        \begin{equation*}
            \begin{bmatrix}
                1 & 2 & 1 & 0 \\
                3 & 9 & 2 & 1 \\
                1 & 2 & 4 & 2 \\
                0 & 1 & 2 & 3
            \end{bmatrix}
            \xrightarrow[\text{stride 1}]{2 \times 2 \text{ max pool}}
            \begin{bmatrix}
                9 & 9 & 2 \\
                9 & 9 & 4 \\
                2 & 4 & 4
            \end{bmatrix}
        \end{equation*}
        Notice something interesting: the maximum value $9$ is repeated four times in the output. Because the original 9 belongs to \textbf{four different overlapping $2 \times 2$ windows}, and it is the maximum in all of them.

        \textcolor{Green3}{\faIcon{question-circle} \textbf{So why would I ever want this?}} Imagine that 9 means that the convolutional filter detected its feature very strongly at that position. After stride-1 max pooling, we're saying that the feature is still present in the \textbf{local neighborhood} of that position. So stride-1 max pooling is a way to \textbf{preserve local maxima} and \textbf{suppress weaker activations within each neighborhood}, rather than aggressively downsampling the feature map.
    \end{itemize}
\end{examplebox}

\begin{figure}[htbp]
    \centering
    \resizetikz{\linewidth}{!}{%
    \tikzsetnextfilename{pooling-stride-1-vs-2-comparison}
    \begin{tikzpicture}[
        >=Stealth,
        font=\small,
        header/.style={font=\normalsize\bfseries, align=center},
        sublabel/.style={font=\footnotesize, text=black!75, align=center},
        badge/.style={draw=#1!80!black, fill=#1!15, text=#1!90!black, font=\scriptsize\bfseries, inner sep=2pt, rounded corners=2pt}
    ]

        % ========================================================
        % LEFT PANEL: STANDARD POOLING (K=2, S=2)
        % ========================================================
        \begin{scope}[shift={(0,0)}]
            % Title & Subtitle centered over left panel (x=2.8)
            \node[header] at (2.8, 4.4) {Standard: $K=2\times 2,\ S=2$};
            \node[sublabel] at (2.8, 3.95) {(Non-Overlapping $\to$ Halved Dimensions)};

            % Input Grid 4x4 (Size: 2.8cm x 2.8cm, cell size = 0.7cm)
            \fill[blue!15]   (0, 1.4)   rectangle (1.4, 2.8);
            \fill[orange!15] (1.4, 1.4) rectangle (2.8, 2.8);
            \fill[teal!15]   (0, 0)     rectangle (1.4, 1.4);
            \fill[purple!15] (1.4, 0)   rectangle (2.8, 1.4);

            \draw[step=0.7cm, black!25] (0,0) grid (2.8,2.8);
            \draw[line width=1.1pt, draw=black!70] (0,0) rectangle (2.8,2.8);
            \draw[line width=0.8pt, draw=black!50, dashed] (1.4,0) -- (1.4,2.8);
            \draw[line width=0.8pt, draw=black!50, dashed] (0,1.4) -- (2.8,1.4);

            % Window Outlines & Badges
            \draw[line width=1.3pt, draw=blue!85!black] (0, 1.4) rectangle (1.4, 2.8);
            \node[badge=blue] at (0.7, 2.1) {Pos 1};

            \draw[line width=1.3pt, draw=orange!85!black] (1.4, 1.4) rectangle (2.8, 2.8);
            \node[badge=orange] at (2.1, 2.1) {Pos 2};

            % Stride jump arrow (S = 2)
            \draw[->, line width=1.1pt, draw=blue!85!black] 
                (0.7, 2.95) to[out=40, in=140] 
                node[above=2pt, font=\scriptsize\bfseries] {Stride $S=2$} (2.1, 2.95);

            % Input Dimension Brace
            \draw[decorate, decoration={brace, amplitude=4pt, mirror}, thick, black!60] 
                (0, -0.2) -- (2.8, -0.2) node[midway, below=5pt, font=\footnotesize] {$N_{\text{in}} = 4$};

            % Mapping Arrow
            \draw[->, line width=1.2pt, black!60] (3.1, 1.4) -- (3.9, 1.4) 
                node[midway, above=2pt, font=\footnotesize\bfseries] {MaxPool};

            % Output Grid 2x2
            \begin{scope}[shift={(4.2, 0.7)}]
                \fill[blue!20]   (0, 0.7)   rectangle (0.7, 1.4);
                \fill[orange!20] (0.7, 0.7) rectangle (1.4, 1.4);
                \fill[teal!20]   (0, 0)     rectangle (0.7, 0.7);
                \fill[purple!20] (0.7, 0)   rectangle (1.4, 0.7);

                \draw[step=0.7cm, black!40] (0,0) grid (1.4,1.4);
                \draw[line width=1.1pt, black!70] (0,0) rectangle (1.4,1.4);

                \node[font=\footnotesize\bfseries, blue!85!black]   at (0.35, 1.05) {$m_1$};
                \node[font=\footnotesize\bfseries, orange!85!black] at (1.05, 1.05) {$m_2$};
                \node[font=\footnotesize\bfseries, teal!85!black]   at (0.35, 0.35) {$m_3$};
                \node[font=\footnotesize\bfseries, purple!85!black] at (1.05, 0.35) {$m_4$};
            \end{scope}

            % Output Dimension Brace
            \draw[decorate, decoration={brace, amplitude=4pt, mirror}, thick, black!60] 
                (4.2, 0.5) -- (5.6, 0.5) node[midway, below=5pt, font=\footnotesize] {$N_{\text{out}} = 2$};

            % Arithmetic / Summary
            \node[align=center, font=\small] at (2.8, -1.5) {
                $\displaystyle N_{\text{out}} = \left\lfloor\frac{4-2}{2}\right\rfloor+1 = 2$\\[4pt]
                \footnotesize Spatial size: $25\%$ ($75\%$ discarded)
            };
        \end{scope}

        % ========================================================
        % VERTICAL SEPARATOR LINE
        % ========================================================
        \draw[black!20, line width=1pt, dashed] (6.1, -2.1) -- (6.1, 4.6);

        % ========================================================
        % RIGHT PANEL: OVERLAPPING POOLING (K=2, S=1)
        % ========================================================
        \begin{scope}[shift={(6.7, 0)}]
            % Title & Subtitle centered over right panel (x=3.15)
            \node[header] at (3.15, 4.4) {Overlapping: $K=2\times 2,\ S=1$};
            \node[sublabel] at (3.15, 3.95) {(Window Step $= 1$ $\to$ Dense Feature Map)};

            % Input Grid 4x4 (Fills first)
            \fill[blue!18] (0, 1.4)   rectangle (1.4, 2.8);
            \fill[red!18]  (0.7, 1.4) rectangle (2.1, 2.8);

            % Grid lines and outer frame
            \draw[step=0.7cm, black!25] (0,0) grid (2.8,2.8);
            \draw[line width=1.1pt, draw=black!70] (0,0) rectangle (2.8,2.8);

            % Window Outlines
            \draw[line width=1.3pt, draw=blue!85!black] (0, 1.4) rectangle (1.4, 2.8);
            \draw[line width=1.3pt, draw=red!80!black, dashed] (0.7, 1.4) rectangle (2.1, 2.8);

            % Enhanced & Non-Colliding Badges
            \node[badge=blue] at (0.35, 2.1) {Pos 1};
            \node[badge=red]  at (1.75, 2.1) {Pos 2};

            % Stride jump arrow (S = 1)
            \draw[->, line width=1.1pt, draw=red!80!black] 
                (0.7, 2.95) to[out=45, in=135] 
                node[above=2pt, font=\scriptsize\bfseries] {Stride $S=1$} (1.4, 2.95);

            % Input Dimension Brace
            \draw[decorate, decoration={brace, amplitude=4pt, mirror}, thick, black!60] 
                (0, -0.2) -- (2.8, -0.2) node[midway, below=5pt, font=\footnotesize] {$N_{\text{in}} = 4$};

            % Mapping Arrow
            \draw[->, line width=1.2pt, black!60] (3.1, 1.4) -- (3.9, 1.4) 
                node[midway, above=2pt, font=\footnotesize\bfseries] {MaxPool};

            % Output Grid 3x3 (Fully Integrated)
            \begin{scope}[shift={(4.2, 0.35)}]
                % Cell Fills
                \fill[blue!20] (0, 1.4)   rectangle (0.7, 2.1);
                \fill[red!20]  (0.7, 1.4) rectangle (1.4, 2.1);

                % Grid Lines over the full 3x3 area
                \draw[step=0.7cm, black!30] (0,0) grid (2.1,2.1);
                \draw[line width=1.1pt, black!70] (0,0) rectangle (2.1,2.1);

                % Matching Cell Borders for m1 and m2
                \draw[line width=1.2pt, draw=blue!85!black] (0, 1.4) rectangle (0.7, 2.1);
                \draw[line width=1.2pt, draw=red!80!black, dashed] (0.7, 1.4) rectangle (1.4, 2.1);

                % Output values
                \node[font=\footnotesize\bfseries, blue!85!black] at (0.35, 1.75) {$m_1$};
                \node[font=\footnotesize\bfseries, red!85!black]  at (1.05, 1.75) {$m_2$};
            \end{scope}

            % Output Dimension Brace
            \draw[decorate, decoration={brace, amplitude=4pt, mirror}, thick, black!60] 
                (4.2, 0.15) -- (6.3, 0.15) node[midway, below=5pt, font=\footnotesize] {$N_{\text{out}} = 3$};

            % Arithmetic / Summary
            \node[align=center, font=\small] at (3.15, -1.5) {
                $\displaystyle N_{\text{out}} = \left\lfloor\frac{4-2}{1}\right\rfloor+1 = 3$\\[4pt]
                \footnotesize Dense local maximum (no halving)
            };
        \end{scope}

    \end{tikzpicture}%
    }
    \caption{Side-by-side comparison of standard non-overlapping pooling ($S=2$) and overlapping pooling ($S=1$).}
    \label{fig:pooling-stride-comparison}
\end{figure}